% =====================================================================
%  Komplex számok
%
%  JAVÍTOTT FORRÁSHIBÁK
%   [p3] z^2 = r^2(cos 2φ + i sin φ)  ->  i sin 2φ
%   [p4] a gyökformula második tagjából hiányzott a sin
%   [p5] φ = -128√3 / -128 = √3  ->  tg φ = √3
%
%  KIEGÉSZÍTÉS
%   [p2] a hányadosnál a "Ha z_2 != 0" feltétel
% =====================================================================

\documentclass[11pt]{article}
\input{preamble}
\input{macros}

\begin{document}
\titleblock{Komplex számok}

% ===================================================================
\section{Algebrai alak}

% [p1]
A valós számok közt az $x^{2}=-1$ egyenletnek nincs megoldása.

\begin{defbox}
\begin{definition}[Képzetes egység]
Jelölje $\ii$ azt a számot, melyre teljesül, hogy
\[
\ii^{2}=-1 .
\]
\end{definition}
\end{defbox}
Ez az $\ii$ egy új, nem valós szám.

\begin{defbox}
\begin{definition}[Komplex szám]
Legyen $a, b\in\R$. A
\[
z=a+b\ii
\]
alakú számokat \emph{komplex számoknak} hívjuk. A komplex számok halmaza: $\C$.
\end{definition}
\end{defbox}

\begin{exbox}
\begin{example}
$z=5+2\ii$.
\end{example}
\end{exbox}

\begin{defbox}
\begin{definition}[Valós rész, képzetes rész, algebrai alak]
A $z=a+b\ii$ komplex számban $a$ a \emph{valós rész} és $b$ a \emph{képzetes rész}. A $z=a+b\ii$ alakban megadott komplex számot a komplex szám \emph{algebrai alakjának} hívjuk.
\end{definition}
\end{defbox}

\begin{exbox}
\begin{example}
$z=4-3\ii$, itt $a=4$ és $b=-3$.
\end{example}
\end{exbox}

\begin{defbox}
\begin{definition}[Konjugált]
A $z=a+b\ii$ komplex szám \emph{konjugáltja}
\[
\bar z=a-b\ii .
\]
\end{definition}
\end{defbox}

\begin{exbox}
\begin{example}
$z=9+2\ii$, $\bar z=9-2\ii$.
\end{example}
\end{exbox}

% ===================================================================
\section{Műveletek komplex számokkal}

% [p1]
Legyen $z_{1}=a_{1}+b_{1}\ii$ és $z_{2}=a_{2}+b_{2}\ii$.

\begin{thmbox}
\begin{proposition}[Összeg, különbség, szorzat algebrai alakban]
\begin{align}
z_{1}+z_{2}&=(a_{1}+a_{2})+(b_{1}+b_{2})\ii, \label{eq:osszeg}\\
z_{1}-z_{2}&=(a_{1}-a_{2})+(b_{1}-b_{2})\ii, \label{eq:kulonbseg}\\
z_{1}z_{2}&=a_{1}a_{2}-b_{1}b_{2}+(a_{1}b_{2}+a_{2}b_{1})\ii. \label{eq:szorzat}
\end{align}
\end{proposition}
\end{thmbox}
\begin{proof}
Az összeg és a különbség a tagok átrendezésével adódik:
\[
z_{1}+z_{2}=(a_{1}+b_{1}\ii)+(a_{2}+b_{2}\ii)=(a_{1}+a_{2})+(b_{1}+b_{2})\ii,
\]
\[
z_{1}-z_{2}=(a_{1}+b_{1}\ii)-(a_{2}+b_{2}\ii)=(a_{1}-a_{2})+(b_{1}-b_{2})\ii .
\]
A szorzatnál a zárójelek felbontása után a $\ii^{2}$ tagot $-1$-gyel helyettesítjük:
\[
z_{1}z_{2}=(a_{1}+b_{1}\ii)(a_{2}+b_{2}\ii)
=a_{1}a_{2}+a_{1}b_{2}\ii+a_{2}b_{1}\ii+b_{1}b_{2}\underbrace{\ii^{2}}_{-1}
=a_{1}a_{2}-b_{1}b_{2}+(a_{1}b_{2}+a_{2}b_{1})\ii .
\]
\end{proof}

% [p2]
\begin{exbox}
\begin{example}
Legyen $z_{1}=4+3\ii$ és $z_{2}=8-5\ii$. Számítsuk ki az összeget, a különbséget és a szorzatot.
\end{example}
\end{exbox}
\begin{solution}
\[
z_{1}+z_{2}=(4+3\ii)+(8-5\ii)=12-2\ii,
\]
\[
z_{1}-z_{2}=(4+3\ii)-(8-5\ii)=-4+8\ii,
\]
\[
z_{1}z_{2}=(4+3\ii)(8-5\ii)=32-20\ii+24\ii-15\underbrace{\ii^{2}}_{-1}=47+4\ii .
\]
\end{solution}

A $\dfrac{z_{1}}{z_{2}}$ hányados algebrai alakját a nevező konjugáltjával való bővítés adja.

\begin{thmbox}
\begin{proposition}[Hányados algebrai alakban]
Ha $z_{2}\neq 0$, akkor
\begin{equation}\label{eq:hanyados}
\frac{z_{1}}{z_{2}}=\frac{a_{1}a_{2}+b_{1}b_{2}}{a_{2}^{2}+b_{2}^{2}}
+\frac{a_{2}b_{1}-a_{1}b_{2}}{a_{2}^{2}+b_{2}^{2}}\,\ii .
\end{equation}
\end{proposition}
\end{thmbox}
\begin{proof}
\begin{align*}
\frac{z_{1}}{z_{2}}&=\frac{a_{1}+b_{1}\ii}{a_{2}+b_{2}\ii}
=\frac{(a_{1}+b_{1}\ii)(a_{2}-b_{2}\ii)}{(a_{2}+b_{2}\ii)(a_{2}-b_{2}\ii)}
=\frac{a_{1}a_{2}-a_{1}b_{2}\ii+a_{2}b_{1}\ii-b_{1}b_{2}\underbrace{\ii^{2}}_{-1}}
{a_{2}^{2}-a_{2}b_{2}\ii+a_{2}b_{2}\ii-b_{2}^{2}\underbrace{\ii^{2}}_{-1}}\\[0.6em]
&=\frac{a_{1}a_{2}+b_{1}b_{2}+(a_{2}b_{1}-a_{1}b_{2})\ii}{a_{2}^{2}+b_{2}^{2}},
\end{align*}
és a valós, illetve a képzetes részt szétírva adódik \eqref{eq:hanyados}.
\end{proof}

\begin{exbox}
\begin{example}
Írjuk algebrai alakba a $\dfrac{4+3\ii}{8-5\ii}$ hányadost.
\end{example}
\end{exbox}
\begin{solution}
\[
\frac{4+3\ii}{8-5\ii}=\frac{(4+3\ii)(8+5\ii)}{(8-5\ii)(8+5\ii)}
=\frac{32+20\ii+24\ii+15\underbrace{\ii^{2}}_{-1}}{64+40\ii-40\ii-25\underbrace{\ii^{2}}_{-1}}
=\frac{17+44\ii}{89}=\frac{17}{89}+\frac{44}{89}\ii .
\]
\end{solution}

% ===================================================================
\section{Trigonometrikus alak}

% [p2]
\begin{defbox}
\begin{definition}[Trigonometrikus alak]
Legyen $z=a+b\ii$, továbbá
\[
a=r\cos\varphi,\qquad b=r\sin\varphi .
\]
Ekkor
\[
z=a+b\ii=r\cos\varphi+r\sin\varphi\cdot \ii=r(\cos\varphi+\ii\sin\varphi),
\qquad 0\le\varphi<2\pi,
\]
és ezt a komplex szám \emph{trigonometrikus alakjának} nevezzük. Itt
\[
r=\sqrt{a^{2}+b^{2}},\qquad \tg\varphi=\frac{b}{a}.
\]
\end{definition}
\end{defbox}

\begin{notefig}{A $z=a+b\ii$ szám trigonometrikus alakjának mennyiségei.}
\begin{tikzpicture}[scale=1.25,>=Stealth]
\draw[->] (-0.4,0) -- (2.9,0);
\draw[->] (0,-0.4) -- (0,2.4);
\coordinate (P) at (2.1,1.75);
\draw[thick,accent] (0,0) -- (P);
\draw[thmaccent,dashed] (P) -- (2.1,0);
\draw[thmaccent,dashed] (P) -- (0,1.75);
\draw[accent] (0.5,0) arc (0:39.8:0.5);
\node[accent,font=\footnotesize] at (0.68,0.2) {$\varphi$};
\node[accent,font=\footnotesize,above left] at (1.1,0.85) {$r$};
\fill (P) circle (1.3pt) node[above right,font=\footnotesize] {$(a, b)$};
\node[font=\footnotesize,below] at (1.05,-0.05) {$r\cos\varphi$};
\node[font=\footnotesize,left] at (-0.05,0.9) {$r\sin\varphi$};
\end{tikzpicture}
\end{notefig}

\begin{exbox}
\begin{example}
Írjuk trigonometrikus alakba a $z=a+b\ii$ alakban adott komplex számot.
\begin{enumerate}[label=\alph*)]
\item $z=2+2\ii$
\item $z=1-\sqrt{3}\,\ii$
\item $z=8\ii$
\end{enumerate}
\end{example}
\end{exbox}
\begin{solution}
a) Itt $a=2$, $b=2$, tehát $r=\sqrt{4+4}=\sqrt{8}$ és
\[
\tg\varphi=\frac{2}{2}=1 \quad\Longrightarrow\quad \varphi=\chosen{\frac{\pi}{4}}\ \text{v.}\ \frac{5\pi}{4},
\]
\[
z=2+2\ii=\sqrt{8}\left(\cos\frac{\pi}{4}+\ii\sin\frac{\pi}{4}\right).
\]

\begin{notefig}{A $z=2+2\ii$ szám az $r$ hosszal és a $\varphi$ szöggel.}
\begin{tikzpicture}[scale=0.9,>=Stealth]
\draw[->] (-0.6,0) -- (1.9,0);
\draw[->] (0,-0.6) -- (0,1.9);
\coordinate (P) at (1.4,1.4);
\draw[thick,accent] (0,0) -- (P);
\draw[thmaccent,dashed] (P) -- (1.4,0);
\draw[thmaccent,dashed] (P) -- (0,1.4);
\draw[accent] (0.45,0) arc (0:45:0.45);
\node[accent,font=\footnotesize] at (0.68,0.2) {$\varphi$};
\node[accent,font=\footnotesize,above left] at (0.75,0.7) {$r$};
\fill (P) circle (1.2pt);
\node[font=\footnotesize,below] at (1.4,-0.08) {$2$};
\node[font=\footnotesize,left] at (-0.08,1.4) {$2$};
\end{tikzpicture}
\end{notefig}

% [p3]
b) Itt $a=1$, $b=-\sqrt{3}$, tehát $r=\sqrt{1+3}=2$ és
\[
\tg\varphi=\frac{-\sqrt{3}}{1}=-\sqrt{3} \quad\Longrightarrow\quad
\varphi=\frac{2\pi}{3}\ \text{v.}\ \chosen{\frac{5\pi}{3}},
\]
\[
z=2\left(\cos\frac{5\pi}{3}+\ii\sin\frac{5\pi}{3}\right).
\]

\begin{notefig}{A $z=1-\sqrt{3}\,\ii$ szám a $\varphi=\frac{5\pi}{3}$ szöggel.}
\begin{tikzpicture}[scale=0.9,>=Stealth]
\draw[->] (-1.6,0) -- (1.9,0);
\draw[->] (0,-2.2) -- (0,1.3);
\coordinate (P) at (1.0,-1.732);
\draw[thick,accent] (0,0) -- (P);
\draw[accent] (0.5,0) arc (0:300:0.5);
\node[accent,font=\footnotesize] at (0.32,0.72) {$\tfrac{5\pi}{3}$};
\fill (P) circle (1.2pt);
\node[font=\footnotesize,above] at (1.0,0.06) {$1$};
\node[font=\footnotesize,left] at (-0.06,-1.732) {$-\sqrt{3}$};
\end{tikzpicture}
\end{notefig}

c) Itt $a=0$, $b=8$, tehát $r=\sqrt{0^{2}+8^{2}}=8$, és
\[
\tg\varphi=\frac{8}{0}
\]
nem értelmezett. Így
\[
z=8\left(\cos\frac{\pi}{2}+\ii\sin\frac{\pi}{2}\right).
\]

\begin{notefig}{A $z=8\ii$ szám a $\varphi=\frac{\pi}{2}$ szöggel.}
\begin{tikzpicture}[scale=0.9,>=Stealth]
\draw[->] (-1.2,0) -- (1.6,0);
\draw[->] (0,-0.6) -- (0,2.0);
\coordinate (P) at (0,1.6);
\draw[thick,accent] (0,0) -- (P);
\draw[accent] (0.45,0) arc (0:90:0.45);
\node[accent,font=\footnotesize] at (0.48,0.55) {$\tfrac{\pi}{2}$};
\fill (P) circle (1.2pt) node[right,font=\footnotesize] {$(0, 8)$};
\end{tikzpicture}
\end{notefig}
\end{solution}

% ===================================================================
\section{Műveletek trigonometrikus alakban adott komplex számokkal}

% [p3]
Legyen
\[
z_{1}=r_{1}(\cos\varphi_{1}+\ii\sin\varphi_{1}),\qquad
z_{2}=r_{2}(\cos\varphi_{2}+\ii\sin\varphi_{2}).
\]

\begin{thmbox}
\begin{proposition}[Szorzat trigonometrikus alakban]
\begin{equation}\label{eq:polszorzat}
z_{1}z_{2}=r_{1}r_{2}\bigl(\cos(\varphi_{1}+\varphi_{2})+\ii\sin(\varphi_{1}+\varphi_{2})\bigr).
\end{equation}
\end{proposition}
\end{thmbox}
\begin{proof}
A zárójelek felbontása és a $\ii^{2}=-1$ helyettesítés után
\[
z_{1}z_{2}=r_{1}(\cos\varphi_{1}+\ii\sin\varphi_{1})\,r_{2}(\cos\varphi_{2}+\ii\sin\varphi_{2})
\]
\[
=r_{1}r_{2}\Bigl(\cos\varphi_{1}\cos\varphi_{2}+\ii\sin\varphi_{1}\cos\varphi_{2}
+\ii\sin\varphi_{2}\cos\varphi_{1}+\underbrace{\ii^{2}}_{-1}\sin\varphi_{1}\sin\varphi_{2}\Bigr)
\]
\[
=r_{1}r_{2}\Bigl(
\underbrace{\cos\varphi_{1}\cos\varphi_{2}-\sin\varphi_{1}\sin\varphi_{2}}_{\cos(\varphi_{1}+\varphi_{2})}
+\ii\bigl(\underbrace{\sin\varphi_{1}\cos\varphi_{2}+\sin\varphi_{2}\cos\varphi_{1}}_{\sin(\varphi_{1}+\varphi_{2})}\bigr)\Bigr)
\]
\[
=r_{1}r_{2}\bigl(\cos(\varphi_{1}+\varphi_{2})+\ii\sin(\varphi_{1}+\varphi_{2})\bigr).
\]
\end{proof}

Így $z=r(\cos\varphi+\ii\sin\varphi)$ esetén
\[
z^{2}=r^{2}(\cos 2\varphi+\ii\sin 2\varphi),
\]
\[
z^{3}=z^{2}\cdot z=r^{2}(\cos 2\varphi+\ii\sin 2\varphi)\,r(\cos\varphi+\ii\sin\varphi)
=r^{3}(\cos 3\varphi+\ii\sin 3\varphi),
\]
\[
\vdots
\]

\begin{thmbox}
\begin{theorem}[Hatványozás pozitív egész kitevőre]
Ha $z=r(\cos\varphi+\ii\sin\varphi)$ és $n\in\Z^{+}$, akkor
\begin{equation}\label{eq:hatvany}
z^{n}=r^{n}(\cos n\varphi+\ii\sin n\varphi).
\end{equation}
\end{theorem}
\end{thmbox}

% [p4]
\begin{thmbox}
\begin{proposition}[Hányados, reciprok, negatív kitevő]
Megmutatható, hogy
\begin{align}
\frac{z_{1}}{z_{2}}&=\frac{r_{1}(\cos\varphi_{1}+\ii\sin\varphi_{1})}{r_{2}(\cos\varphi_{2}+\ii\sin\varphi_{2})}
=\frac{r_{1}}{r_{2}}\bigl(\cos(\varphi_{1}-\varphi_{2})+\ii\sin(\varphi_{1}-\varphi_{2})\bigr),
\label{eq:polhanyados}\\
\frac{1}{z}&=\frac{1(\cos 0+\ii\sin 0)}{r(\cos\varphi+\ii\sin\varphi)}
=\frac{1}{r}\bigl(\cos(-\varphi)+\ii\sin(-\varphi)\bigr),
\label{eq:reciprok}
\end{align}
és általában $z=r(\cos\varphi+\ii\sin\varphi)$ esetén
\begin{equation}\label{eq:negativkitevo}
z^{-n}=r^{-n}\bigl(\cos(-n\varphi)+\ii\sin(-n\varphi)\bigr).
\end{equation}
\end{proposition}
\end{thmbox}

\begin{thmbox}
\begin{theorem}[Hatványozás egész kitevőre]
Összefoglalva $z=r(\cos\varphi+\ii\sin\varphi)$ esetén, ha $m\in\Z$, akkor
\begin{equation}\label{eq:egeszhatvany}
z^{m}=r^{m}(\cos m\varphi+\ii\sin m\varphi).
\end{equation}
\end{theorem}
\end{thmbox}

% ===================================================================
\section{Gyökvonás}

% [p4]
\begin{defbox}
\begin{definition}[$n$-edik gyök]
A $z\in\C$ komplex számnak a $w\in\C$ szám \emph{$n$-edik gyöke}, ha
\[
w^{n}=z .
\]
\end{definition}
\end{defbox}

\begin{thmbox}
\begin{theorem}[A komplex $n$-edik gyökök]
A $z\neq 0$ komplex számnak $n$ db.\ $n$-edik gyöke van, és
\begin{equation}\label{eq:gyok}
z^{1/n}=r^{1/n}\left(\cos\frac{\varphi+k2\pi}{n}+\ii\sin\frac{\varphi+k2\pi}{n}\right),
\qquad k=0, 1, \dots, n-1 .
\end{equation}
\end{theorem}
\end{thmbox}
\begin{proof}
Legyen
\[
z=r(\cos\varphi+\ii\sin\varphi),\qquad w=\rho(\cos\alpha+\ii\sin\alpha).
\]
Így
\[
r(\cos\varphi+\ii\sin\varphi)=z=w^{n}=\rho^{n}(\cos n\alpha+\ii\sin n\alpha),
\]
ahonnan
\[
r=\rho^{n}\ \Longrightarrow\ \rho=r^{1/n},
\]
\[
n\alpha=\varphi+k\cdot 2\pi\ \Longrightarrow\ \alpha=\frac{\varphi+k2\pi}{n},\qquad k\in\Z .
\]
Csak $k=0, 1, \dots, n-1$ esetén kapunk különböző komplex számokat, így a $z\neq 0$ komplex számnak $n$ db.\ $n$-edik gyöke van.
\end{proof}

% [p5]
\begin{exbox}
\begin{example}
Adjuk meg algebrai alakban a $z=-128-128\sqrt{3}\,\ii$ szám negyedik gyökeit.
\end{example}
\end{exbox}
\begin{solution}
Itt $a=-128$ és $b=-128\sqrt{3}$, tehát
\[
r=\sqrt{(-128)^{2}+(-128\sqrt{3})^{2}}=256,
\]
\[
\tg\varphi=\frac{-128\sqrt{3}}{-128}=\sqrt{3}\quad\Longrightarrow\quad
\varphi=\frac{\pi}{3}\ \text{v.}\ \chosen{\frac{4\pi}{3}} .
\]

\begin{notefig}{A $z=-128-128\sqrt{3}\,\ii$ szám a $\varphi=\frac{4\pi}{3}$ szöggel.}
\begin{tikzpicture}[scale=0.9,>=Stealth]
\draw[->] (-2.0,0) -- (1.6,0);
\draw[->] (0,-2.4) -- (0,1.4);
\coordinate (P) at (-1.0,-1.732);
\draw[thick,accent] (0,0) -- (P);
\draw[accent] (0.5,0) arc (0:240:0.5);
\node[accent,font=\footnotesize] at (0.32,0.72) {$\tfrac{4\pi}{3}$};
\fill (P) circle (1.2pt);
\node[font=\footnotesize,above] at (-1.0,0.06) {$-128$};
\node[font=\footnotesize,left] at (-0.06,-1.9) {$-128\sqrt{3}$};
\end{tikzpicture}
\end{notefig}

\noindent A \eqref{eq:gyok} formula szerint
\[
z^{1/4}=256^{1/4}\left(\cos\frac{\frac{4\pi}{3}+k2\pi}{4}+\ii\sin\frac{\frac{4\pi}{3}+k2\pi}{4}\right),
\qquad k=0, 1, 2, 3 .
\]
Az egyes esetek:
\begin{align*}
k=0:\quad z_{1}&=256^{1/4}\left(\cos\frac{\pi}{3}+\ii\sin\frac{\pi}{3}\right)
=4\left(\frac{1}{2}+\ii\,\frac{\sqrt{3}}{2}\right)=2+2\sqrt{3}\,\ii,\\
k=1:\quad z_{2}&=256^{1/4}\left(\cos\frac{5\pi}{6}+\ii\sin\frac{5\pi}{6}\right)
=4\left(-\frac{\sqrt{3}}{2}+\ii\cdot\frac{1}{2}\right)=-2\sqrt{3}+2\ii,\\
k=2:\quad z_{3}&=256^{1/4}\left(\cos\frac{4\pi}{3}+\ii\sin\frac{4\pi}{3}\right)
=4\left(-\frac{1}{2}+\ii\left(-\frac{\sqrt{3}}{2}\right)\right)=-2-2\sqrt{3}\,\ii,\\
k=3:\quad z_{4}&=256^{1/4}\left(\cos\frac{11\pi}{6}+\ii\sin\frac{11\pi}{6}\right)
=4\left(\frac{\sqrt{3}}{2}+\ii\left(-\frac{1}{2}\right)\right)=2\sqrt{3}-2\ii .
\end{align*}
\end{solution}

\begin{exbox}
\begin{example}
Oldjuk meg a
\[
z^{2}+4z+13=0
\]
egyenletet.
\end{example}
\end{exbox}
\begin{solution}
A megoldóformulából
\[
z_{1,2}=\frac{-4\pm\sqrt{16-52}}{2}=\frac{-4\pm\sqrt{-36}}{2}.
\]
A $\sqrt{-36}$ meghatározásához vegyük észre, hogy itt $a=-36$ és $b=0$, tehát
\[
-36=36(\cos\pi+\ii\sin\pi),
\]
\[
\sqrt{-36}=36^{1/2}\left(\cos\frac{\pi+k2\pi}{2}+\ii\sin\frac{\pi+k2\pi}{2}\right),\qquad k=0, 1 .
\]

\begin{notefig}{A $-36$ szám a $\varphi=\pi$ szöggel.}
\begin{tikzpicture}[scale=0.9,>=Stealth]
\draw[->] (-2.0,0) -- (1.6,0);
\draw[->] (0,-1.0) -- (0,1.4);
\coordinate (P) at (-1.4,0);
\draw[thick,accent] (0,0) -- (P);
\draw[accent] (0.5,0) arc (0:180:0.5);
\node[accent,font=\footnotesize] at (0.1,0.68) {$\pi$};
\fill (P) circle (1.2pt);
\node[font=\footnotesize,below] at (-1.4,-0.08) {$-36$};
\end{tikzpicture}
\end{notefig}

\noindent A két eset:
\[
k=0:\quad 36^{1/2}\left(\cos\frac{\pi}{2}+\ii\sin\frac{\pi}{2}\right)=6\ii,
\]
\[
k=1:\quad 36^{1/2}\left(\cos\frac{3\pi}{2}+\ii\sin\frac{3\pi}{2}\right)=-6\ii .
\]
Ezekkel
\[
z_{1,2}=\frac{-4\pm(\pm 6\ii)}{2}=\frac{-4\pm 6\ii}{2}=
\begin{cases}
-2+3\ii,\\
-2-3\ii .
\end{cases}
\]
\end{solution}

\end{document}
